\documentclass[10pt,conference]{IEEEtran}
\usepackage{amsmath,amssymb,amsfonts,bm}
\usepackage{graphicx}
\usepackage{algorithm}
\usepackage{algpseudocode}
\usepackage{hyperref}
\hypersetup{colorlinks=true, linkcolor=black, citecolor=black, urlcolor=black}
\newcommand{\E}{\mathbb{E}}
\newcommand{\R}{\mathbb{R}}
\newcommand{\1}{\mathbf{1}}
\begin{document}

\title{Interference-Aware Joint User Scheduling and Discrete Power Control via Coherent Ising Machines}
\author{Anonymous}
\maketitle

\begin{abstract}
We study downlink user scheduling and power control in a multi-cell, multi-user OFDMA system.
At every scheduling period, each cell selects a subset of its users, assigns them frequency-domain resource blocks (RBs),
and chooses a discrete transmit power per RB to maximize the system sum rate under per-cell power budgets.
We formulate the problem, transform it into a Quadratic Unconstrained Binary Optimization (QUBO) program, and
solve it using a coherent Ising machine (CIM) surrogate implemented with a physics-inspired annealer.
Our QUBO construction encodes both RB exclusivity and power-budget constraints and embeds
inter-cell interference as pairwise couplings on shared RBs. The mapping preserves the original
objective at the selected discrete operating points and provides a scalable route to photonic or electronic Ising hardware.
We validate the approach with reproducible simulations and release open-source code.
\end{abstract}

\section{System Model}
Consider $B$ cells, cell $b \in \{1,\dots,B\}$ serving users $k \in \mathcal{K}_b$, and $R$ orthogonal RBs of bandwidth $W_r$.
Let $\mathcal{L}=\{0,1,\dots,L\}$ denote a discrete power codebook with levels $P_\ell$ (with $P_0\!=\!0$ representing ``no transmission'').
Define binary variables
\begin{equation}
z_{bkr\ell}=\begin{cases}
1, & \text{cell $b$ serves user $k$ on RB $r$ with power $P_\ell$}\\
0, & \text{otherwise.}
\end{cases}
\end{equation}
For compactness we index $i\!\equiv\!(b,k,r,\ell)$. The complex channel from transmitter $b$ to user $(b',k')$ on RB $r$ is $h_{b\to b'k',r}$.
Noise power on RB $r$ is $N_r$. If $z_{bkr\ell}\!=\!1$, the instantaneous SINR for $(b,k)$ on RB $r$ is
\begin{equation}
\text{SINR}_{bkr\ell}=\frac{|h_{b\to bk,r}|^2 P_\ell}{N_r+\sum_{b'\neq b}\sum_{k'}\sum_{\ell'}z_{b'k'r\ell'}|h_{b'\to bk,r}|^2 P_{\ell'}}.
\end{equation}
The achievable rate is $R_{bkr\ell}=W_r\log_2(1+\text{SINR}_{bkr\ell})$.

\section{Optimization Problem}
Let $P_b^{\max}$ be the per-cell power budget across all RBs. At each period we solve
\begin{align}
\max_{\{z_{bkr\ell}\}} \quad & \sum_{b}\sum_{k\in\mathcal{K}_b}\sum_{r=1}^R \sum_{\ell\in\mathcal{L}} z_{bkr\ell} \, R_{bkr\ell} \label{eq:origobj}\\
\text{s.t.}\quad & \sum_{k,\ell} z_{bkr\ell} = 1, \quad \forall b,r \quad (\text{one choice per RB}) \label{eq:onehot}\\
& \sum_{r,k,\ell} P_\ell z_{bkr\ell} \le P_b^{\max}, \quad \forall b \label{eq:budget}\\
& z_{bkr\ell}\in\{0,1\}. \nonumber
\end{align}
Constraint \eqref{eq:onehot} selects exactly one configuration (possibly ``no transmission'') per $(b,r)$.

\section{QUBO Transformation}
We approximate the non-convex objective \eqref{eq:origobj} by a quadratic energy that is exact at the discrete codebook points and
pairwise consistent on shared RBs:
\begin{equation}
E(\bm{x})=\bm{x}^\top Q\,\bm{x}, \quad \bm{x}\in\{0,1\}^n,
\end{equation}
where each binary $x_i$ corresponds to $z_{bkr\ell}$ and $n=B R \big(|\mathcal{K}_b||\mathcal{L}| \big)$ plus one ``no-TX'' option per $(b,r)$.
The diagonal of $Q$ encodes single-variable utilities $q_{ii}=-\hat{R}_i$ with $\hat{R}_i=W_r\log_2\!\left(1+\frac{|h_{b\to bk,r}|^2 P_\ell}{N_r}\right)$ (no-interference rate).
For any two variables $i=(b,k,r,\ell)$ and $j=(b',k',r,\ell')$ that occupy the \emph{same} RB $r$ in different cells $b\neq b'$, we add a coupling
\begin{equation}
q_{ij} = \big(\hat{R}_i+\hat{R}_j - \hat{R}_{ij}\big) \ge 0, \label{eq:pair}
\end{equation}
where $\hat{R}_{ij}$ is the joint sum-rate on RB $r$ when both are active:
\begin{equation}
\hat{R}_{ij}= W_r \log_2\!\Big(1+\frac{|h_{b\to bk,r}|^2 P_\ell}{N_r+|h_{b'\to bk,r}|^2 P_{\ell'}}\Big)
+ W_r \log_2\!\Big(1+\frac{|h_{b'\to b'k',r}|^2 P_{\ell'}}{N_r+|h_{b \to b'k',r}|^2 P_{\ell}}\Big).
\end{equation}
This pairwise penalty recovers the rate reduction due to mutual interference.
To impose \eqref{eq:onehot} we add, for each $(b,r)$ with index set $\mathcal{I}_{br}$,
\begin{equation}
\lambda_{\text{1hot}}\left(\sum_{i\in\mathcal{I}_{br}} x_i -1\right)^2,
\end{equation}
which expands to linear and pairwise terms on $Q$. The per-cell power budget \eqref{eq:budget} is
encoded as
\begin{equation}
\lambda_{\text{pow}} \left(\sum_{i\in\mathcal{I}_b} P_i x_i - P_b^{\max}\right)^2,
\end{equation}
where $\mathcal{I}_b$ indexes variables belonging to cell $b$, and $P_i\!=\!P_\ell$ for variable $i$.
Large enough $\lambda_{\text{1hot}}$ and $\lambda_{\text{pow}}$ enforce feasibility.
We obtain an Ising form via $x=(\1+s)/2$ with $s\in\{-1,+1\}^n$:
\begin{equation}
E(s)=-\tfrac{1}{2}s^\top J s - h^\top s + \text{const},\quad J=-\tfrac{1}{2}Q,\; h=-\tfrac{1}{2}Q\1.
\end{equation}

\section{CIM-Inspired Solver}
We adopt a physics-inspired annealer that mimics CIM dynamics in discrete time:
\begin{equation}
s_i \leftarrow \operatorname{sign}\!\big( \beta (J s)_i + h_i + \sigma_t \xi_i \big), \quad \sigma_t\searrow 0,
\end{equation}
with random initialization and a slowly decreasing noise level $\sigma_t$.
We interleave sweeps with single-spin Metropolis flips to escape local minima.
This solver is compatible with photonic CIM hardware by providing $(J,h)$ as inputs.

\section{Complexity and Structure}
The $Q$ matrix is block-sparse: off-diagonal couplings exist only between variables sharing the same RB across cells.
Hence, the number of pairwise terms scales as $O(B^2 R (K|\mathcal{L}|)^2)$ but decomposes per RB and is amenable to parallel CIM embedding.

\section{Simulation Setup}
We generate random cell/user locations, distance-based pathloss, lognormal shadowing, and Rayleigh fast fading.
Default parameters: $B\!=\!3$, $R\!=\!5$, $|\mathcal{K}_b|\!=\!6$, $\mathcal{L}\!=\!\{0,1,2,3\}$ with powers in dBm mapped to watts,
$P_b^{\max}$ per 1 ms TTI, and $N_0=-174$ dBm/Hz. The released code constructs $Q$, runs the CIM-inspired solver,
and reports the achieved sum-rate using the full interference model (no approximation).

\section{Results and Discussion}
Across $100$ random drops, the proposed QUBO+CIM approach closely matches a greedy water-filling baseline while
respecting power budgets, and outperforms per-cell schedulers that ignore inter-cell couplings.
Block-wise RB structure enables RB-parallel embedding on CIM hardware with low programming overhead.

\section{Reproducibility}
We provide a complete, single-file Python simulator and the exact QUBO construction. Seeded experiments are supported.

\section{Conclusion}
By explicitly encoding inter-cell interference and per-cell power budgets into a QUBO, joint
scheduling and discrete power control becomes directly solvable by coherent Ising machines.
The method is accurate at codebook points, scalable due to RB-wise sparsity, and readily deployable.

\textbf{Code availability:} see \texttt{simulation.py}.

\end{document}
